\begin{problem}[理论考试][太阳内部核反应]{32}
  太阳是我们赖以生存的恒星。它的主要成分是氢，氢在自身引力的作用下收缩而导致升温，当温度高到一定程度时，中性原子将电离成质子和电子组成的等离子体，并在其核心区域达到约 $1.5\times10^{7}K$ 的高温和 $1.6\times10^{5}kg/m^{3}$ 以上的高密度，产生热核聚变而放出巨大能量，从而抗衡自身的引力收缩达到平衡，而成为恒星。太阳内部主要核反应过程为
  \begin{equation}
    {}^{1}\mathrm{H} + {}^{1}\mathrm{H} \rightarrow \mathrm{D} + \mathrm{e}^{+} + v_{e}
    \label{eq:I.32}
  \end{equation}
  \begin{equation}
    \mathrm{D} + {}^{1}\mathrm{H} \rightarrow {}^{3}\mathrm{He} + \mathrm{x}
    \label{eq:II.32}
  \end{equation}
  \begin{equation}
    {}^{3}\mathrm{He} + {}^{3}\mathrm{He} \rightarrow {}^{4}\mathrm{He} + {}^{1}\mathrm{H} + {}^{1}\mathrm{H}
    \label{eq:III.32}
  \end{equation}
  其中第一个反应的概率由弱相互作用主导，概率很低，这恰好可以使得能量缓慢释放。反应产物正电子 $e^{+}$ 会与电子 $e^{-}$ 湮灭为 $\gamma$ 射线，即
  \begin{equation}
    \mathrm{e}^{+} + \mathrm{e}^{-} \rightarrow \gamma + \gamma
    \label{eq:IV.32}
  \end{equation}
  已知：质子($^{1}$H)、氘(D)、氦-3($^{3}$He)、氦-4($^{4}$He)和电子的质量分别为938.27、1875.61、2808.38、3727.36和0.51（MeV/$c^{2}$）（误差为0.01 MeV/$c^{2}$），c为真空中的光速，中微子 $v_{e}$ 的质量小于3 eV/$c^{2}$。普朗克常量 $h=6.626\times10^{-34}$ J·s，$c=3.0\times10^{8}$ m/s，波尔兹曼常量 $k=1.381\times10^{-23}$ J/K，电子电量 $e=1.602\times10^{-19}$ C。

  \begin{subproblem}
    试用理想气体模型估算处于热平衡状态的各种粒子的平均动能及太阳核心区的压强（请分别用 eV 和 atm 为单位）；
  \end{subproblem}
  \begin{subproblem}
    反应式（II）中的 x 是什么粒子（$\alpha$、$\beta$、$\gamma$、p 和 n 之一）？请计算该粒子的动能和动量的大小，是否可以找到一个参照系，使得 x 粒子的动能为零？
  \end{subproblem}
  \begin{subproblem}
    给出反应式（I）中各反应产物的动能的范围；
  \end{subproblem}
  \begin{subproblem}
    反应式（IV）中的 $\gamma$ 射线（光子）经过很多次的康普顿散射到达太阳表面的时候，其波长约为 $5.4 \times 10^{-7} \, m$。设入射光子的能量为 $E_{0}$，一次散射后光子的能量为 E，光子的散射角为 $\phi$（出射光子动量相对于入射光子动量的夹角）。推导出由 $E_{0}$ 和 $\phi$ 表示的 E 的表达式；假设从太阳中心到表面，光子经历了 $10^{26}$ 次康普顿散射，且每次散射的角度相同，求每一次散射的散射角 $\phi$。
  \end{subproblem}
\end{problem}